\documentclass[../my_ntust_thesis.tex]{subfiles}
\begin{document}
\chapter{緒論}
\label{sec:intro}

\section{背景與功能分割演進}

\textcolor{blue}{為了滿足 5G 網路多樣化且嚴苛的需求，第三代合作夥伴計畫 (3GPP) 探索了各種功能分割 (Functional Split) 選項，將傳統單體式基地台架構進行解構 \cite{3gpp_tr_38_801}。在此基礎上，O-RAN 聯盟進一步將解構後的架構標準化為三個主要邏輯節點：開放式中央單元 (O-CU)、開放式分散單元 (O-DU) 以及開放式無線單元 (O-RU)。這種解構架構允許營運商根據其具體的部署環境，優化傳輸頻寬與計算資源。}

\textcolor{blue}{目前業界公認有三種主流的功能分割配置。第一種是 \textbf{Split Option 2}，其定義了 O-CU 與 O-DU 之間的介面。第二種是 \textbf{Split Option 7.2}，這是 O-RAN 標準化中，位於 O-DU 與 O-RU 之間的低層分割 (Low-layer Split) \cite{oran_wg4_cus}。第三種是 \textbf{Split Option 6}，即 MAC-PHY 分割（通常稱為 O-DU 高低層分割），這也是本論文研究的核心。為了確保 Option 6 的互操作性，小型基地台論壇 (Small Cell Forum, SCF) 已將介面規範統一於網路功能應用平台介面 (nFAPI) 之下 \cite{SCF225}。}

\textcolor{blue}{這三種主流分割選項在同步需求與管理機制上存在根本差異：}

\begin{itemize}
\item \textcolor{blue}{\textbf{Split Option 2 (CU/DU 分割)：} 分割點位於 PDCP 與 RLC 層之間。由於其運作於較高的協定層，其同步需求是基於 \textbf{序號 (Sequence Numbers)} 而非嚴格的時間域。核心管理機制依賴 \textbf{流量控制與深層緩衝 (Flow control and Deep Buffering)} 來處理封包傳輸。因此，它非常適合部署於 \textbf{一般封包交換網路}，不需精確的物理層定時對齊。}
\item \textcolor{blue}{\textbf{Split Option 6 (DU 高低層分割)：} 運作於 MAC 與 PHY 層之間，此分割需要更嚴格的 \textbf{訊框/時槽 (Frame/Slot)} 級別同步。其核心管理高度聚焦於 \textbf{延遲管理 (Delay management)}，以確保排程指令在其預定的時槽邊界前送達。與 Option 2 類似，它旨在透過 \textbf{一般封包交換網路} 運作，但必須自適應地調整以應對網路抖動。}
\item \textcolor{blue}{\textbf{Split Option 7.2 (DU/RU 分割)：} 位於物理層內部，此分割要求極高的 \textbf{相位/符號 (Phase/Symbol)} 級別同步。其核心管理需要嚴格的硬體輔助機制，例如 \textbf{硬體 PTP 與時窗 (Window)} 強制執行。由於其對延遲與抖動高度敏感，因此必須部署於專門的 \textbf{時效性網路 (Time-Sensitive Networking, TSN)} 環境。}
\end{itemize}

\textcolor{blue}{表~\ref{tab:split_comparison} 總結了這三種分割架構在同步維護方面的技術特性。}

\begin{table}[h]
\centering
\caption{功能分割與定時維護之比較}
\label{tab:split_comparison}
\small
\begin{tabularx}{\textwidth}{@{}L L L L@{}}
\toprule
\textbf{分割選項} & \textbf{同步層級} & \textbf{核心管理} & \textbf{部署環境} \\ \midrule
\textbf{Option 2} (CU/DU 分割) & 序號 & 流量控制與深層緩衝 & 一般封包交換網路 \\
\textbf{Option 6} (DU 高低層分割) & 訊框/時槽 & 延遲管理 & 一般封包交換網路 \\
\textbf{Option 7.2} (DU/RU 分割) & 相位/符號 & 硬體 PTP 與時窗 & 時效性網路 (TSN) \\ \bottomrule
\end{tabularx}
\end{table}

\textcolor{blue}{這種架構差異意味著當面對非理想的前傳網路 (Fronthaul) 時，Option 6 相比於 Option 7.2 具有更好的延遲容忍度，同時與 Option 2 相比則提供了更細粒度的物理層控制。}

\section{研究動機與問題定義}

\textcolor{blue}{O-RAN 前傳工作小組規範 \cite{oran_wg4_cus_v02} 強調了研究低層分割 (LLS) 架構中高延遲問題的必要性。具體而言，該規範的附錄 L 討論了在非理想傳輸環境下運作的考量因素。長延遲的前傳鏈路可能會破壞傳統的混合自動重傳請求 (HARQ) 時序，僅為基地台留下緊迫的 4~ms 處理預算。這種限制通常需要使用即時作業系統 (RTOS) 與專用硬體。然而，對於中低影響的部署場景，如室內小型基地台與固定無線存取 (FWA)，HARQ 延遲或預測性 HARQ 的性能影響是可以忽略不計的 \cite{oran_wg4_cus_v02}。因此，強制執行亞毫秒級的延遲約束反而限制了創新且具成本效益的部署探索。}

\textcolor{blue}{在此背景下，nFAPI (Split Option 6) 架構天生符合非理想前傳部署的需求。透過將 MAC 與 PHY 層分離，Split 6 與 Split 7.2x 相比展現出更高的延遲容忍度（通常為 2~ms 至 10~ms）。這種增強的容忍度使得託管 MAC 層的虛擬網路功能 (VNF) 能夠部署在雲端環境中的通用處理器 (GPP) 上，消除了對昂貴專用光纖與精確同步硬體的嚴格要求。此外，nFAPI 內建了定時程序，使 MAC 層能夠根據前傳延遲狀況調整其排程行為，直接解決了 HARQ 時窗失效的問題。}

\textcolor{blue}{將 MAC 層遷移到集中化雲端環境的原理在於資源池化與彈性。在 5G vRAN 架構中，物理層 (L1) 的計算負載主要由相對靜態的參數決定，例如系統頻寬與天線配置。相比之下，第二層 (L2) 協定疊的處理複雜度，特別是 MAC 排程器，與活躍使用者設備 (UE) 的數量及流量體積呈線性甚至二次方的相依關係 \cite{Foukas2021_ICC, ComplexityPFS2019_WCNC}。例如，隨著連接 UE 數量的增加，MAC 排程器必須持續評估比例公平指標、處理緩衝狀態報告並管理邏輯通道，導致 CPU 週期大幅增加 \cite{GPP_CSCN2018}。透過將 MAC 層解構為 VNF，雲端基礎設施可以動態分配虛擬核心以應對波動的使用者密度，與本地端 PNF 部署相比，能維持具成本效益的擴展。}

\textcolor{blue}{然而，這種資源擴展為同步帶來了新的挑戰。隨著 UE 數量的增加，在給定的傳輸時間間隔 (TTI) 內，MAC 處理時間 ($T_{\text{process\_VNF}}$) 不可避免地會延長。這段延長的處理時間壓縮了前傳傳輸延遲可用的剩餘時間裕度。如果排程機制沒有考慮到這種內生性增加的處理延遲，靜態的定時配置將導致排程訊息過晚到達 PNF。儘管 Option 6 具有靈活的定時調整能力，但其動態調整演算法在現有的開源實作（如 OpenAirInterface, OAI）中仍未成熟。因此，當面臨非確定性的作業系統排程、與 UE 相關的處理負載以及網路擁塞時，系統經常會遇到「Missing Subframe」錯誤。本論文選擇 Split 6 作為主要研究對象，提出了一種「自適應提前時間控制 (Adaptive Advance-Time Control)」架構，用以動態補償處理與傳輸延遲的變動，從而有效地克服這些同步失效問題。關於這種非確定性延遲的詳細分析以及相關問題的驗證將在第~\ref{sec:experiment} 章中呈現，特別是在情境 I (\ref{subsec:scenario_1}) 中對這些問題進行了分離討論。}

\section{相關研究}
\textcolor{blue}{無線存取網路 (RAN) 架構從分散式 RAN (D-RAN) 到開放式 RAN (O-RAN) 的演進，需要對功能分割進行深入研究 \cite{larsen2019comparison, 10186379, 11124199}。在非理想傳輸網路上部署解構式 RAN 的一個基本挑戰是管理前傳延遲與功能分割限制的影響。O-RAN 架構引入了多種功能分割以適應不同的網路環境，每種分割都提出了不同的同步需求。例如，最近關於 O-RAN 中無細胞大規模 MIMO (CF-mMIMO) 的研究分析了功能分割 7.2x 與 8 下的混合前傳規劃，強調了前傳容量與處理位置如何影響密集部署中的總體成本與性能 \cite{Mohammed2026Fronthaul}。在這些基於 7.2x 的架構中，嚴格的前傳定時要求意味著任何因抖動而超出 RU 定時時窗的封包都會被被動丟棄（即時窗外封包丟棄），使得系統對非理想傳輸條件高度敏感。相比之下，我們的設計採用了基於 nFAPI 的 Option 6 MAC--PHY 分割，它能夠在 MAC--PHY 邊界進行微秒級的主動相位與定時調整。這使得系統能夠透過主動對齊傳輸時間來部分吸收前傳抖動，而不是像傳統 Option 7.2x 前傳處理鏈那樣被動地丟棄延遲封包。}

\textcolor{blue}{目前的研究廣泛聚焦於優化虛擬化與雲端 RAN 中的資源分配與延遲管理。如表~\ref{tab:related_works_comparison} 所述，頂尖會議的近期研究從各個架構層級解決了延遲問題。例如，Adamuz-Hinojosa \textit{et al.} \cite{Adamuz2024ORANUS} 提出了 ORANUS，這是一個針對 6G O-RAN 量身定制的延遲編排框架。雖然他們的工作利用隨機網路演算在 10~ms 到 1~s 的時間尺度上有效地界定了超高可靠低延遲通訊 (URLLC) 佇列，但其運作於較高的編排層級，而非解決微秒級的同步問題。}

\textcolor{blue}{此外，將 MAC 排程器虛擬化會引入顯著的伺服器處理變動。Lv \textit{et al.} \cite{Lv2024QoS} 透過聯合優化排程與編排來解決此問題，以確保 5G vRAN 的服務品質 (QoS)。他們的研究強調，將 MAC 處理轉移到雲端 VNF 伺服器會引入內生性的處理延遲，進而壓縮了允許的前傳傳輸延遲裕度。然而，他們的優化主要針對計算資源競爭，而非在像 Split Option 6 這樣嚴格的時槽級同步語境下動態適應傳輸網路抖動。}

\textcolor{blue}{Chen \textit{et al.} \cite{chen2022m3} 進一步證明了在解構式架構中進行極其嚴格、微秒級定時控制的必要性，他們為 C-RAN 架構下的多細胞 MIMO 網路設計了一個亞毫秒級排程器。他們的工作證明了亞毫秒排程的嚴苛限制，說明靜態定時提前 (Timing Advance) 配置（常規用於 OpenAirInterface 等開源框架中）在網路抖動下極易產生時框缺失 (Missing Subframe) 錯誤。這突顯了我們研究所填補的關鍵空白：雖然現有的頂尖文獻聚焦於高層編排或集中式基頻處理，本論文則專門解決 Option 6 固有的微秒相位級同步挑戰。我們提出了一種動態的雙迴路自適應提前時間控制演算法，並在端到端商用測試平台（商用 UE 與 Pegatron O-RU）上進行了驗證，證明了其應對非確定性前傳抖動與系統變動的有效性。}

\begin{table}[h]
\centering
\caption{相關研究之比較分析}
\label{tab:related_works_comparison}
\footnotesize
\begin{tabularx}{\textwidth}{@{}L L L L L L@{}}
\toprule
\textbf{文獻} & \textbf{架構與分割} & \textbf{目標延遲問題} & \textbf{時間尺度} & \textbf{核心機制} & \textbf{評估平台} \\ \midrule
\cite{Mohammed2026Fronthaul} TWC 2026 & O-RAN (FS 7.2x \& 8) & 前傳容量與成本權衡 & 靜態規劃 & 網路優化 & 數值模擬 \\
\cite{Adamuz2024ORANUS} INFOCOM 2024 & O-RAN (基於 RIC) & URLLC 佇列與傳輸延遲 & 10ms - 1s & 透過 SNC 的資源編排 & Python 系統模擬 \\
\cite{chen2022m3} INFOCOM 2022 & C-RAN (BBU-RRH) & 集中化基頻處理延遲 & $<$ 1ms & GPU 並行化聯合排程器 & 商用 GPU + 模擬 \\
\cite{Lv2024QoS} INFOCOM 2024 & 5G vRAN (CU-DU) & 伺服器非確定性計算 & 跨時間尺度 & 聯合排程與編排 & OAI 原型 + USRP \\
\textbf{本論文} & \textbf{O-RAN Option 6 (nFAPI)} & \textbf{非確定性前傳抖動} & \textbf{微秒 (相位級)} & \textbf{自適應提前時間控制} & \textbf{OAI E2E 商用測試平台} \\ \bottomrule
\end{tabularx}
\end{table}

\section{論文貢獻}
\begin{itemize}
\item \textcolor{blue}{\textbf{nFAPI Split 6 的延遲管理：} 我們分析並量化了非確定性的延遲因素，包括網路抖動與伺服器處理時間，這些因素導致了解構式 RAN 部署中的同步失效。}
\item \textcolor{blue}{\textbf{自適應提前時間控制演算法：} 我們提出了一種利用雙迴路架構的動態定時調整演算法。「收斂優化迴路 (Convergence Optimization Loop)」估計端到端抖動，而「VNF 定時致動器 (VNF Timing Actuator)」則調整提前時槽 ($S_{ahead}$)，以確保 P7 訊息的準時送達。}
\item \textcolor{blue}{\textbf{實驗驗證：} 我們在 OAI 平台上實作了所提出的機制。實驗結果顯示，我們的演算法消除了「Missing Subframe」錯誤，並在非理想前傳環境下維持了穩定的吞吐量。}
\end{itemize}
\end{document}
